%% kapa: exam=1 round=22 year=2011 period=2 form=A subject=accounting subject_ko=회계학 seq=1 number=1
%% source: Q-Net 감정평가사 기출문제 (cst00309) artlSeq=5036790 file=extracted/1008/5036790/2011년 제22회 감정평가사 1차시험 기출문제/2교시.hwp
%% conversion: kordoc 4.24.51 -> tools/gen1.py (auto); review=proofread-gpt-astra; flags=none
\begin{question}{1}
자산의 측정 및 평가에 관한 설명으로 옳지 않은 것은?
\begin{choices}
  \item 역사적원가로 자산을 평가하는 경우 미실현이익을 인식하지 않는다.
  \item 현행원가로 자산을 평가하는 경우 실물자본유지개념에 적합하다.
  \item 현행유출가치는 청산가치의 측정척도가 된다.
  \item 자산을 공정가치로 평가하면 역사적원가로 평가하는 경우보다 회계이익이 과대계상될 수 있다.
  \item 현금흐름할인가치는 자산의 정의에 가장 충실한 측정기준이지만 실질가치를 반영하지 못한다.
\end{choices}
\end{question}
\answer{5}
